\section*{Bab \ref{ch:b2:plant-meristems} --- Perkembangan Vegetatif Tumbuhan: Meristem dan Pertumbuhan}

\begin{solution}{exo:b2:plant-meristems:1}
Sebuah \omterm{def:b2:plant-meristems:meristem}{meristem} adalah populasi sel yang tak terdiferensiasi dan
membelah yang bertahan, dan yang darinya jaringan tumbuhannya berasal.
\omterm{def:b2:plant-meristems:meristem}{Meristem apikal pucuk}: daun, batang, \omterm{def:b2:angiosperm-reproduction:flower}{bunga}; \omterm{def:b2:plant-meristems:meristem}{meristem apikal akar}:
akarnya; \omterm{def:b2:plant-meristems:wood}{kambium pembuluh}: xilem sekunder (\omterm{def:b2:plant-meristems:wood}{kayu}) dan floem; \omterm{def:b2:plant-meristems:meristem}{kambium}
gabus: \omterm{def:b2:plant-meristems:wood}{periderm} (gabus) \omterm{def:b2:plant-meristems:wood}{kulit kayunya}.
\end{solution}

\begin{solution}{exo:b2:plant-meristems:2}
\omterm{prop:b2:plant-meristems:ram}{Tudung akar} (perlindungan, pelumasan, penginderaan gravitasi, selnya
terkelupas); \omterm{def:b2:plant-meristems:meristem}{meristem} beserta \omterm{prop:b2:plant-meristems:ram}{pusat diamnya} (pembelahan); \omterm{prop:b2:plant-meristems:ram}{zona
pemanjangan} (selnya memanjang sepuluh sampai dua puluh kali lipat);
zona \omterm{def:b2:cell-differentiation:differentiation}{diferensiasi} (\omterm{prop:b2:plant-meristems:ram}{rambut akar}, xilem dan floem yang matang). Akar
sampingnya muncul dari perisikel akar yang matang, berseberangan dengan
kutub xilemnya.
\end{solution}

\begin{solution}{exo:b2:plant-meristems:3}
Tumbuhan buncis dipenggal: tunas sampingnya tumbuh keluar. Pada
perangkat kedua, tunggul potongannya menerima sebongkah agar berisi
auksin: tunasnya tetap dorman; pada perangkat kendalinya, sebongkah
agar polos: tunasnya tumbuh. Kesimpulannya: auksin dari pucuk
apikalnyalah yang menghambat tunas sampingnya --- \omterm{prop:b2:plant-meristems:apical}{dominansi apikal}
bersifat hormon.
\end{solution}

\begin{solution}{exo:b2:plant-meristems:4}
Berumur enam puluh tahun; \omterm{def:b2:plant-meristems:wood}{kayu} gubalnya adalah lima belas lingkaran
terakhirnya; yang hidup adalah \omterm{def:b2:plant-meristems:meristem}{kambiumnya}, floemnya, sel jari-jari dan
parenkima \omterm{def:b2:plant-meristems:wood}{kayu} gubalnya, serta \omterm{def:b2:plant-meristems:meristem}{kambium} gabusnya --- sebuah selaput
tipis mengelilingi sebuah teras yang mati.
\end{solution}

\begin{solution}{exo:b2:plant-meristems:5}
$0.1\times(P - 0.3)$: \qty{0.01}{h^{-1}} (berlipat dua dalam
$\ln 2/0.01 =
\qty{69}{h}$), \qty{0.03}{h^{-1}} (\qty{23}{h}),
\qty{0.05}{h^{-1}} (\qty{14}{h}). Pada \qty{0.25}{MPa} turgornya berada
di bawah ambang luluhnya: pertumbuhannya berhenti, sebelum ada layu.
\end{solution}

\begin{solution}{exo:b2:plant-meristems:6}
Daun 1--8 pada 0, 137,5, 275, 52,5, 190, 327,5, 105, dan
\qty{242.5}{\degree}. Diurutkan: 0, 52,5, 105, 137,5, 190, 242,5, 275,
327,5 --- celah terkecilnya \qty{32.5}{\degree}. Tiap daunnya duduk di
sebuah celah daun sebelumnya, sehingga kedelapannya saling menaungi
sesedikit mungkin dan tajuknya mengisi cakram cahayanya secara merata.
\end{solution}

\begin{solution}{exo:b2:plant-meristems:7}
$2\pi r h\Delta r$: pada \qty{5}{cm}, $2\pi\times 0.05\times 12\times
0.003 = \qty{0.011}{m^3}$, \qty{5.7}{kg}, \qty{2.8}{kg} karbon; pada
\qty{40}{cm}, \qty{0.090}{m^3}, \qty{45}{kg}, \qty{23}{kg} karbon
--- delapan kali lebih banyak dari lebar lingkaran yang sama.
\end{solution}

\begin{solution}{exo:b2:plant-meristems:8}
\emph{clv3}: tanpa rem --- \omterm{def:b2:plant-meristems:meristem}{meristem} yang luar biasa besar, daun dan
organ \omterm{def:b2:angiosperm-reproduction:flower}{bunga} tambahan, batang yang berfasiasi. \emph{wus}: tanpa pedal
gas --- \omterm{def:b2:plant-meristems:meristem}{meristemnya} kehabisan tenaga sesudah beberapa daun lalu
didirikan ulang berulang-ulang. Mutan gandanya: \omterm{def:b2:meiosis-heredity:vocabulary}{fenotipe} \emph{wus},
sehingga WUS bekerja di hilir CLV3: tugas CLV3 adalah menekan WUS, dan
tanpa WUS tidak ada yang dapat ditekan.
\end{solution}

\begin{solution}{exo:b2:plant-meristems:9}
Pengangkutannya: \qty{20}{h}. Difusinya: $(0.2)^{2}/(2\times 10^{-9}) =
\qty{2e7}{s}$, yaitu kira-kira 230 hari. Pengangkutan terkutub membuat
sebuah isyarat hormon dapat dipakai sepanjang tumbuhannya dalam sehari,
dan memberinya sebuah arah.
\end{solution}

\begin{solution}{exo:b2:plant-meristems:10}
\qty{3}{cm} jaringan matang sehari dengan sel sepanjang
\qty{200}{\micro m}: \num{150} sel tiap berkasnya tiap hari. Maka
seratus sel yang membelah harus masing-masing membelah 1,5 kali sehari:
sebuah daur kira-kira \qty{16}{h}.
\end{solution}

\begin{solution}{exo:b2:plant-meristems:11}
Giberelin memanjangkan ruas batangnya; tumbuhan yang tidak dapat
menanggapinya memiliki ruas yang pendek dan batang yang pendek. Pupuk
nitrogen yang berat membuat gandum yang tinggi tumbuh lebih tinggi dan
bermalai lebih berat sampai ia rebah lalu membusuk; batang yang pendek
dan kaku tetap berdiri, dan menaruh karbon tambahannya ke dalam
bulirnya alih-alih ke dalam jeraminya, sehingga indeks panennya naik.
Di alam liar, tumbuhan yang pendek dilampaui dan dinaungi tetangganya
lalu kalah.
\end{solution}

\begin{solution}{exo:b2:plant-meristems:12}
\omterm{def:b2:plant-meristems:meristem}{Meristem} membiarkan sebatang tumbuhan menambahkan organ di tempat dan
pada waktu keadaannya memihaknya --- daun ke arah cahaya, akar ke arah
air --- lalu meninggalkan organ yang tidak berbayar; tiap tunas
ketiaknya adalah sebuah titik tumbuh cadangan, sehingga penggembalaan,
kebakaran, dan pemangkasan dilalui lewat pertumbuhan kembali; dan
karena \omterm{def:b2:plant-meristems:meristem}{meristemnya} embrionik selamanya, sebuah \omterm{def:b2:asexual-reproduction:clone}{genet} tidak memiliki
umur yang tetap. Ongkosnya: sebatang tumbuhan tidak dapat berpindah
maupun menata ulang apa yang telah dibangunnya --- sebuah organ yang
telah ditempatkan tinggal di tempatnya, dan sebuah tubuh yang dibangun
lewat penimbunan harus mengusung masa lalunya yang mati (\omterm{def:b2:plant-meristems:wood}{kayu} terasnya)
bersamanya.
\end{solution}

\begin{solution}{pb:b2:plant-meristems:1}
\textbf{1.} $150/3 = 50$ daun tiap pucuk apikalnya; \num{10000} bagi
pohonnya.
\textbf{2.} 0, 137,5, 275, 52,5, dan \qty{190}{\degree}.
\textbf{3.} $8\times 137.5 = 1100 = 3\times 360 + 20$; $13\times 137.5
= 1787.5 = 5\times 360 - 12.5$; $21\times 137.5 = 2887.5 = 8\times 360
+ 7.5$: daun ke-22 adalah yang pertama berada dalam selang
\qty{10}{\degree} dari daun pertamanya.
\textbf{4.} Tiap daun barunya duduk di celah yang terlebar, sehingga
daunnya saling menaungi paling sedikit; auksinlah, yang pengangkutan
terkutubnya (pengangkut PIN) menguras sekeliling tiap primordiumnya,
yang menetapkan jaraknya.
\textbf{5.} $\num{10000}\times 80\,\text{cm}^{2} = \qty{80}{m^2}$:
seluruh tajuknya adalah daun musim ini.
\textbf{6.} Seluruhnya; daun pertamanya dibangun dari pati yang
tersimpan di dalam \omterm{def:b2:plant-meristems:wood}{kayu} dan akarnya pada tahun sebelumnya.
\textbf{7.} Pucuk apikal yang terlipat tiga membuat primordium lebih
cepat --- dua atau tiga tiap plastokronnya --- sehingga daun dan luas
daunnya dapat terlipat tiga, tetapi pucuknya akan berfasiasi, dengan
batang seperti pita dan filotaksis yang kacau, dan sebagian besar daun
tambahannya akan menaungi dirinya sendiri.
\textbf{8.} Siang: $0.05\times(0.55 - 0.35) = \qty{0.01}{h^{-1}}$;
malam: $0.05\times 0.40 = \qty{0.02}{h^{-1}}$.
\textbf{9.} Siang: $2e^{0.12} = \qty{2.25}{cm}$, bertambah
\qty{0.25}{cm}; malam: $2.25e^{0.24} = \qty{2.87}{cm}$, bertambah
\qty{0.61}{cm}.
\textbf{10.} Laju puratanya \qty{0.015}{h^{-1}} sepanjang \qty{96}{h}:
$2e^{1.44}
= \qty{8.4}{cm}$.
\textbf{11.} Siang: $0.05\times 0.30 = \qty{0.015}{h^{-1}}$; malam:
$0.05\times 0.50 = \qty{0.025}{h^{-1}}$.
\textbf{12.} $0.05\times 0.05 = \qty{0.0025}{h^{-1}}$, yaitu seperempat
laju siang yang normal. Penyesuaian osmosis menaikkan kandungan zat
terlarut selnya sehingga, pada potensial air yang sama, turgornya naik
kembali di atas ambangnya.
\textbf{13.} Pada siang hari, transpirasinya menurunkan potensial air
daunnya beserta turgornya, sehingga $P - Y$ menjadi kecil; pada malam
hari turgornya pulih. Pertumbuhannya dibatasi air, bukan gula, jam demi
jam.
\textbf{14.} $2\pi\times 0.10\times 8\times 0.005 = \qty{0.025}{m^3}$;
\qty{12.6}{kg} kering.
\textbf{15.} \qty{6.3}{kg} karbon.
\textbf{16.} $80\times 4\times 150 = \qty{48}{kg}$ karbon.
\textbf{17.} $6.3/48 = \qty{13}{\%}$. Lubuk yang lain: daunnya sendiri,
akar, cabang, dan rantingnya, respirasi tiap selnya yang hidup
(kira-kira separuh dari yang difiksasi), cadangannya, \omterm{def:b2:angiosperm-reproduction:flower}{bunga}, dan
\omterm{def:b2:angiosperm-reproduction:seed}{bijinya}.
\textbf{18.} Jari-jari \qty{0.20}{m}: $2\pi\times 0.20\times 8\times
0.005 = \qty{0.050}{m^3}$, dua kali lipat tahun ini --- keliling yang
dilalui \omterm{def:b2:plant-meristems:meristem}{kambiumnya} telah berlipat dua (dan pohonnya akan lebih tinggi).
\textbf{19.} Sebuah tahun yang kering atau dingin; atau penggundulan
daun oleh serangga, embun beku, atau kebakaran. Penyebab iklim
menyempitkan lingkaran yang sama pada tiap pohon di daerahnya lalu
berkorelasi dengan catatan cuacanya; kerusakan mengenai satu pohon atau
satu rumpun, dan embun beku atau kebakaran meninggalkan parut anatomi
di dalam lingkarannya.
\textbf{20.} Tunas yang terdekat dengan tiap potongannya tumbuh keluar
dalam beberapa hari menjadi beberapa pucuk; dalam sebulan pohonnya
menjadi lebih rimbun dan lebih rapat, dengan pucuk yang lebih banyak
tetapi lebih pendek.
\textbf{21.} Auksin pada potongannya menggantikan isyarat pucuk
apikalnya: tunasnya tetap dorman dan tidak ada percabangan yang
menyusul.
\textbf{22.} Tiap pemangkasannya membuang pucuk apikalnya lalu
melepaskan tunas di bawahnya; enam kali setahun, tiap pelepasannya
menambahkan satu susun pucuk pendek, dan permukaannya terisi olehnya.
\textbf{23.} Dari tunas ketiak yang dorman pada tunggulnya dan dari
tunas liar yang dibentuk \omterm{def:b2:plant-meristems:meristem}{kambiumnya}; tiap tunasnya adalah sebuah
\omterm{def:b2:plant-meristems:meristem}{meristem} lengkap yang sanggup membangun kembali sistem pucuknya,
sedangkan kepala seekor hewan adalah sebuah \omterm{def:b2:vertebrate-development:induction}{pengatur} tunggal yang tak
tergantikan dan tanpa salinan cadangan.
\textbf{24.} Sitokinin yang lebih banyak mencapai tunasnya memacu
pertumbuhannya --- menyiram tumbuhan yang dipangkas membuatnya lebih
bercabang.
\textbf{25.} 50 daun tiap pucuk apikal, \num{10000} tiap pohonnya; laju
ruasnya \qty{0.01}{h^{-1}} pada siang hari dan \qty{0.02}{h^{-1}} pada
malam hari; \omterm{def:b2:plant-meristems:wood}{kayu} setahunnya \qty{0.025}{m^3}, \qty{12.6}{kg} kering,
\qty{6.3}{kg} karbon.
\end{solution}
